% from ega4-2-fr.tex
\oldpage[IV-2]{207}

Dans la suite de ce numéro, nous allons appliquer (7.5.1) à l'étude des
produits tensoriels complétés $A\widehat\otimes_k B$ d'anneaux locaux
noethériens qui sont des algèbres sur un corps $k$.

\begin{lemma}[7.5.5]
\label{IV.2.s7.5.5-fr}
Soient $k$ un corps, $A$, $B$ deux anneaux locaux noethériens complets
contenant $k$, le corps résiduel de $A$ étant une extension finie de $k$. Soit
$C$ le produit tensoriel complété $A\widehat\otimes_k B$
(0$_{\mathrm I}$, 7.7.5). Alors :
\begin{enumerate}
\item[\emph{(i)}] $C$ est un anneau semi-local noethérien complet.
\item[\emph{(ii)}] $C$ est un $A$-module plat et un $B$-module plat.
\item[\emph{(iii)}] Si $\mathfrak m$ est l'idéal maximal de $A$,
$\mathfrak mC$ est contenu dans le radical de $C$, et $C/\mathfrak mC$ est
isomorphe à $(A/\mathfrak m)\otimes_k B$.
\item[\emph{(iv)}] Les corps résiduels des composants locaux de $C$ sont des
extensions finies du corps résiduel de $B$.
\end{enumerate}
Les propriétés (i), (iii) et la première assertion de (ii) sont des cas
particuliers de (0, 19.7.1.2). Pour prouver que $C$ est un $B$-module plat,
notons que pour tout $h>0$,
\[
C/\mathfrak m^hC=(A/\mathfrak m^h)\otimes_k B
\]
est un $B$-module plat, puisque $k$ est un corps; il suffit donc d'appliquer
(0$_{\mathrm{III}}$, 10.2.6) à chacun des composants locaux de $C$ (dont $C$
est le composé direct). Enfin, si $\mathfrak n$ est l'idéal maximal de $B$, les
corps résiduels des composants locaux de $C$ sont aussi ceux des composants
locaux de l'anneau artinien
$(A/\mathfrak m)\otimes_k(B/\mathfrak n)$, qui, par hypothèse, sont des
extensions finies de $B/\mathfrak n$.
\end{lemma}

\begin{proposition}[7.5.6]
\label{IV.2.s7.5.6-fr}
Soient $k$ un corps, $A$, $B$ deux anneaux locaux noethériens complets
contenant $k$, et dont les corps résiduels sont des extensions finies de $k$.
Soit d'autre part $\boldsymbol R$ une propriété vérifiant les conditions
$(\mathrm R'_{\mathrm I})$, $(\mathrm R_{\mathrm{III}})$ et
$(\mathrm R_{\mathrm{IV}})$ (la propriété $\boldsymbol P$ étant définie à
partir de $\boldsymbol R$ par (7.5.0)). Supposons que $\boldsymbol R(A)$ soit
vraie, ainsi que $\boldsymbol P(\operatorname{Spec}(B),k)$ (ce qui signifie que
pour toute extension finie $k'$ de $k$ et pour chacun des anneaux locaux
complets $B_i$ dont le composé direct est $B\otimes_k k'$,
$\boldsymbol R(B_i)$ est vraie). Alors, pour chacun des anneaux locaux complets
$C_j$ dont $A\widehat\otimes_k B$ est le composé direct,
$\boldsymbol R(C_j)$ est vraie.
\end{proposition}

\begin{proof}
Il résulte de (7.5.5, (iii)) que chacun des homomorphismes $A\to C_j$ est local
et de (7.5.5, (ii)) que chacun des $C_j$ est un $A$-module plat; enfin, par
(7.5.5, (iv)), les corps résiduels des $C_j$ sont des extensions finies de
$A/\mathfrak m$. D'autre part, la propriété
\[
\boldsymbol P(\operatorname{Spec}(C\otimes_A(A/\mathfrak m)),A/\mathfrak m)
\]
est vraie, car $C/\mathfrak mC=(A/\mathfrak m)\otimes_kB$, et l'assertion
résulte de la définition de $\boldsymbol P$ et de l'hypothèse que
$A/\mathfrak m$ est une extension finie de $k$. Toutes les conditions de
(7.5.1) sont donc vérifiées pour les homomorphismes locaux $A\to C_j$ et les
morphismes correspondants
$\operatorname{Spec}(C_j)\to\operatorname{Spec}(A)$ sont donc des
$\boldsymbol P$-morphismes; la conclusion résulte alors de
$(\mathrm R'_{\mathrm I})$ et de l'hypothèse que $\boldsymbol R(A)$ est vraie.
\end{proof}

\begin{corollary}[7.5.7]
\label{IV.2.s7.5.7-fr}
(Chevalley). --- Soient $k$ un corps parfait (resp. algébriquement clos), $A$,
$B$ deux anneaux locaux noethériens complets contenant $k$ et dont les corps
résiduels sont des extensions finies de $k$. Alors, si $A$ et $B$ sont réduits
(resp. intègres) le produit tensoriel complété $A\widehat\otimes_k B$ est
réduit (resp. intègre).
\end{corollary}

\begin{proof}
I) Supposons $k$ parfait, $A$ et $B$ réduits; soient $A_i$ (resp. $B_j$) les
quotients de $A$ (resp. $B$) par ses idéaux premiers minimaux
($1\leq i\leq r$, $1\leq j\leq s$), qui sont complets; l'hypothèse entraîne
que $A$ (resp. $B$) s'identifie à un sous-anneau du composé direct des $A_i$
(resp. $B_j$); les produits tensoriels étant pris sur un corps,
$A\otimes_kB$ s'identifie à un sous-anneau du composé direct des
$A_i\otimes_kB_j$, et on vérifie aussitôt que la topologie produit tensoriel
de $A\otimes_kB$ est induite par le produit des topologies des
$A_i\otimes_kB_j$. Il en résulte que $A\widehat\otimes_kB$ s'identifie à un
sous-anneau du composé direct des $A_i\widehat\otimes_kB_j$, et on est donc
ramené au cas où $A$ et $B$ sont intègres. Soient alors $A'$ et $B'$ les
clôtures intégrales de $A$ et $B$ respectivement; on sait, par le théorème de
finitude de Nagata (0, 23.1.6), que $A'$ (resp. $B'$) est un $A$-module (resp.
$B$-module) de type fini et un anneau local complet. En outre,
$A\widehat\otimes_kB$ s'identifie à un sous-anneau de
$A'\widehat\otimes_kB'$ : en effet, il suffira de prouver que
$A\widehat\otimes_kB$ s'identifie à un sous-anneau de
$A'\widehat\otimes_kB$, et ce dernier à un sous-anneau de
$A'\widehat\otimes_kB'$. Cela résulte du lemme suivant :

\begin{lemma}[7.5.7.1]
\label{IV.2.s7.5.7.1-fr}
Soient $A$, $B$ deux anneaux locaux noethériens complets contenant un corps
$k$ et dont les corps résiduels sont des extensions finies de $k$. Soit
$\mathfrak m$ l'idéal maximal de $A$. Pour tout $A$-module $M$ de type fini
(muni de la topologie $\mathfrak m$-adique), le produit tensoriel complété
$M\widehat\otimes_kB$ (0$_{\mathrm I}$, 7.7.1) s'identifie canoniquement à
$M\otimes_A(A\widehat\otimes_kB)$.
\end{lemma}

\begin{proof}
En effet, on a un isomorphisme canonique
$M\otimes_kB\xrightarrow{\sim}M\otimes_A(A\otimes_kB)$ donc un homomorphisme
canonique composé
\[
\varphi:M\otimes_kB\xrightarrow{\sim}M\otimes_A(A\otimes_kB)
\longrightarrow M\otimes_A(A\widehat\otimes_kB)
\]
et il est immédiat que cet homomorphisme est continu pour les topologies
produits tensoriels; en outre, $M\otimes_A(A\widehat\otimes_kB)$ est séparé et
complet ((7.5.5) et (0$_{\mathrm I}$, 7.7.8)), donc on a aussi par complétion
un homomorphisme continu
\[
\widehat\varphi:M\widehat\otimes_kB\longrightarrow
M\otimes_A(A\widehat\otimes_kB).
\]

